\documentclass[
    language=english,
    doctype=book-chapter,
    institution=none,
]{../../omnilatex}

\RequirePackage{config/document-settings}
\addbibresource{../../bib/bibliography.bib}

\begin{document}

\section{Introduction}
This chapter provides an overview of distributed consensus algorithms
and their applications in modern computing systems.  We examine the
fundamental challenges and survey the key approaches developed over
the past four decades.

\section{Theoretical Foundations}
The study of distributed consensus begins with the FLP impossibility
result, which establishes that no deterministic protocol can guarantee
consensus in an asynchronous system with even a single faulty
process~\cite{bibid}.

\subsection{The CAP Theorem}
Brewer's CAP theorem and its formal proof by Gilbert and Lynch establish
that a distributed system cannot simultaneously provide consistency,
availability, and partition tolerance.

\subsection{Consistency Models}
We distinguish between strong consistency (linearizability), sequential
consistency, and eventual consistency, each providing different
guarantees about the order of operations.

\section{Consensus Protocols}

\subsection{Paxos}
Lamport's Paxos protocol provides a foundational approach to consensus
in asynchronous systems.  The protocol operates in two phases: a
prepare phase and an accept phase.

\subsection{Raft}
The Raft protocol was designed as an understandable alternative to Paxos,
providing equivalent safety guarantees with a more straightforward
leader-based approach.

\section{Practical Applications}
Modern distributed systems rely on consensus for critical operations
including configuration management, leader election, and state
replication.

\section{Summary}
This chapter has surveyed the landscape of distributed consensus,
from theoretical foundations to practical implementations.

\printbibliography

\end{document}
